% TEMPLATE-MILC-v3.tex - plantilla maestra para todas las guias MILC v3
% Ver scripts/generadores/build-guia-milc.py para la lista completa de
% claves esperadas. El script valida que no queden placeholders sin
% reemplazar antes de escribir el archivo final.

\documentclass[11pt,letterpaper]{article}
\usepackage[margin=1.75cm]{geometry}
\usepackage[table]{xcolor}
\usepackage{fontspec}
\usepackage[spanish]{babel}
\usepackage{tikz}
\usepackage[most]{tcolorbox}
\usepackage{tabularx}
\usepackage{array}
\usepackage{fancyhdr}
\usepackage{titlesec}
\usepackage{ragged2e}
\usepackage{enumitem}
\usepackage{microtype}

\setmainfont{Helvetica Neue}
\setsansfont{Helvetica Neue}
\setlength{\parindent}{0pt}
\setlength{\parskip}{6pt}
\hyphenpenalty=200
\exhyphenpenalty=200
\pretolerance=400
\tolerance=2000
\setlength{\emergencystretch}{1em}
\renewcommand{\arraystretch}{1.25}
\setlist{leftmargin=5mm,itemsep=1mm,topsep=1mm}
\newcolumntype{Y}{>{\RaggedRight\arraybackslash}X}

\definecolor{milcMagenta}{HTML}{E6007E}
\definecolor{milcVino}{HTML}{5A0038}
\definecolor{milcTurquesa}{HTML}{0093A5}
\definecolor{milcVerde}{HTML}{58A923}
\definecolor{milcMostaza}{HTML}{E5B400}
\definecolor{milcOcre}{HTML}{B86600}
\definecolor{milcNegro}{HTML}{241816}
\definecolor{milcGris}{HTML}{F7F7F4}
\definecolor{milcLinea}{HTML}{E8E2DD}
\definecolor{milcGradePrimary}{HTML}{0066FF}
\definecolor{milcGradeDark}{HTML}{003D99}
\definecolor{milcGradeSoft}{HTML}{D6E8FF}
\definecolor{milcGradeAccent}{HTML}{CFE7FF}
\definecolor{milcGradeText}{HTML}{FFFFFF}
\color{milcNegro}

\pagestyle{fancy}
\fancyhf{}
\lhead{\small Tecnología e Informática · Grado 6}
\rhead{\small Guía 25 · MILC v3}
\cfoot{\small\thepage}
\renewcommand{\headrulewidth}{0pt}

\newtcolorbox{titlebox}[1]{
  enhanced,colback=#1,colframe=#1,arc=7mm,boxrule=0pt,
  left=7mm,right=7mm,top=3mm,bottom=3mm,
  before skip=8pt,after skip=8pt,
  fontupper=\sffamily\bfseries\Large\color{white}
}

\newtcolorbox{softbox}[3]{
  enhanced,colback=#2,colframe=#1,borderline west={4pt}{0pt}{#1},
  arc=2mm,boxrule=.4pt,left=4mm,right=4mm,top=3mm,bottom=3mm,
  before skip=6pt,after skip=8pt,title={#3},coltitle=white,
  colbacktitle=#1,fonttitle=\sffamily\bfseries,fontupper=\RaggedRight,
  attach boxed title to top left={xshift=3mm,yshift=-2mm},
  boxed title style={arc=2mm, boxrule=0pt}
}

\newtcolorbox{drawbox}[2]{
  enhanced,colback=white,colframe=#1,arc=4mm,boxrule=1pt,
  left=5mm,right=5mm,top=4mm,bottom=4mm,before skip=8pt,after skip=8pt,
  title={#2},coltitle=white,colbacktitle=#1,fonttitle=\sffamily\bfseries,
  fontupper=\RaggedRight,
  attach boxed title to top left={xshift=4mm,yshift=-2mm},
  boxed title style={arc=3mm, boxrule=0pt}
}

\newtcolorbox{infoband}[1]{
  enhanced,colback=#1!8,colframe=#1!50,arc=2mm,boxrule=.5pt,
  left=4mm,right=4mm,top=2mm,bottom=2mm,before skip=4pt,after skip=4pt,
  fontupper=\small\RaggedRight
}

% Puente narrativo entre secciones (contrato editorial Conéctate).
% Da continuidad: "Acabas de X. Ahora vas a Y porque Z."
% Visualmente: banda izquierda en color del grado, texto en itálica pequeña.
\newtcolorbox{puentebox}{
  enhanced,colback=milcGradeSoft,colframe=milcGradeDark,
  borderline west={3pt}{0pt}{milcGradeDark},
  arc=0mm,boxrule=0pt,left=5mm,right=4mm,top=2.5mm,bottom=2.5mm,
  before skip=4pt,after skip=8pt,
  fontupper=\itshape\small\color{milcNegro}\RaggedRight
}
\newcommand{\puente}[1]{%
  \begin{puentebox}%
    \textcolor{milcGradeDark}{\textbf{\textrightarrow}}\hspace{2mm}#1%
  \end{puentebox}%
}

\newcommand{\linead}[1]{\noindent\color{milcLinea}\rule{#1}{0.5pt}\color{milcNegro}}
\newcommand{\respuesta}[1]{\par\vspace{2mm}\foreach \n in {1,...,#1}{\noindent\color{milcLinea}\rule{\linewidth}{0.5pt}\par\vspace{5mm}}\color{milcNegro}}
\newcommand{\checkbox}{\textcolor{milcGradeDark}{\fbox{\phantom{x}}}\hspace{5pt}}

\newcommand{\phasecard}[4]{
\begin{tcolorbox}[enhanced,colback=#3,colframe=#2,arc=3mm,boxrule=.5pt,height=3.05cm,valign=center,halign=center,left=1.5mm,right=1.5mm,top=2mm,bottom=2mm]
{\sffamily\bfseries\fontsize{22}{24}\selectfont\textcolor{#2}{#1}}\par
{\sffamily\bfseries\small #4}
\end{tcolorbox}
}

\begin{document}
\thispagestyle{empty}
\begin{tikzpicture}[remember picture,overlay]
  \fill[white] (current page.south west) rectangle (current page.north east);
  \path[fill=milcGradePrimary,rounded corners=16mm]
    ([xshift=.55cm,yshift=-.55cm]current page.north west) rectangle
    ([xshift=-.55cm,yshift=3.65cm]current page.south east);

  \node[anchor=north west,text=milcGradeText] at ([xshift=2.0cm,yshift=-1.85cm]current page.north west)
    {\sffamily\bfseries\fontsize{14}{16}\selectfont GRADO};
  \node[anchor=north west,text=milcGradeText,opacity=.82] at ([xshift=2.0cm,yshift=-2.75cm]current page.north west)
    {\sffamily\bfseries\fontsize{9}{11}\selectfont SERIE GUÍAS · TECNOLOGÍA E INFORMÁTICA};

  \begin{scope}[shift={([xshift=-3.05cm,yshift=-2.55cm]current page.north east)}]
    \fill[white,opacity=.80] (-.62,-.52) rectangle (.62,.52);
    \draw[milcGradeText,line width=1pt] (-.62,-.52) rectangle (.62,.52);
    \fill[milcGradeDark] (-.42,-.38) rectangle (-.22,.06);
    \fill[milcGradeAccent] (-.10,-.38) rectangle (.10,.24);
    \fill[milcGradePrimary!60!black] (.22,-.38) rectangle (.42,.38);
  \end{scope}

  \node[anchor=west,text=milcGradeText] at ([xshift=1.9cm,yshift=-12.85cm]current page.north west)
    {\sffamily\bfseries\fontsize{124}{124}\selectfont 6};
  \draw[milcGradeText,line width=1.7pt]
    ([xshift=4.52cm,yshift=-11.68cm]current page.north west) circle (.13cm);

  \node[anchor=west,text=milcGradeText,text width=8.7cm,align=left] at ([xshift=10.35cm,yshift=-11.6cm]current page.north west)
    {
     {\sffamily\bfseries\fontsize{19}{22}\selectfont Tablas e\\imágenes ---\\mostrar datos\\y agregar\\voz visual}\\[4mm]
     {\sffamily\bfseries\fontsize{23}{26}\selectfont Guía 25-6-TIC}\\[2mm]
     {\sffamily\fontsize{13}{16}\selectfont Período 3 · Procesador de texto e Internet}\\[5mm]
     {\sffamily\bfseries\fontsize{11}{13}\selectfont Institución Educativa Sor María Juliana}\\[1mm]
     {\sffamily\fontsize{10}{12}\selectfont Autor: Dr. Álvaro Cárdenas Orozco}
    };

  \path[fill=milcGradeSoft,rounded corners=7mm]
    ([xshift=1.35cm,yshift=.85cm]current page.south west) rectangle
    ([xshift=-1.35cm,yshift=4.25cm]current page.south east);
  \draw[milcGradeDark,line width=.65pt,rounded corners=7mm]
    ([xshift=1.35cm,yshift=.85cm]current page.south west) rectangle
    ([xshift=-1.35cm,yshift=4.25cm]current page.south east);
  \node[anchor=center,text=milcNegro,text width=17.0cm,align=center] at ([yshift=2.55cm]current page.south)
    {\sffamily\fontsize{10}{12}\selectfont Metodología MILC · Apertura ancestral · Pensamiento computacional · Triángulo Dussel-Estoicismo-Floridi\\
    \textcolor{milcGradeDark}{\bfseries Producto:} Un \textbf{documento} con: \textbf{(a) una tabla} de 5 filas × 4 columnas que muestre datos reales (ejemplo: \emph{horario semanal del colegio} con materia, día, hora, salón); \textbf{(b) 2 imágenes} insertadas con \emph{ajuste de texto alrededor} (una a la izquierda, otra a la derecha), cada una con \emph{pie de imagen} debajo; \textbf{(c) 2 párrafos de texto} relacionados con las imágenes que fluyen alrededor de ellas. Más \textbf{en el cuaderno}: dibujo del documento + 5 reglas de uso de imágenes que aprendiste.};
\end{tikzpicture}
\mbox{}
\newpage

\section*{Apertura: Tablas e imágenes --- mostrar datos y agregar voz visual}

\begin{softbox}{milcGradeDark}{milcGradeSoft}{Saber ancestral}
\textbf{Doña Mercedes la maestra rural de la vereda La Plata de Cartago usaba el tablero de varias maneras al tiempo.} A veces escribía texto. A veces hacía \textbf{cuadros} con líneas para comparar (\emph{``Las 4 estaciones: primavera, verano, otoño, invierno''} en una tabla con sus características). Y a veces pegaba \textbf{recortes de periódico o dibujos} con líneas que apuntaban a las palabras importantes. \emph{``Las palabras solas no siempre alcanzan''}, decía doña Mercedes. \emph{``A veces hay que mostrar la cosa, no solo nombrarla.''} Un dibujo de la araña en el cuaderno enseñaba más que 5 párrafos describiéndola. Una tabla comparando \emph{vacas vs caballos} se entendía en 30 segundos, mientras 2 párrafos comparándolos se demoraban 5 minutos. \textbf{Mezclar texto, tablas e imágenes era ya el saber de doña Mercedes hace 40 años.} Hoy los procesadores de texto te dan exactamente esas mismas herramientas en digital.
\end{softbox}

\begin{softbox}{milcTurquesa}{milcGris}{Saber contemporáneo}
Un documento profesional rara vez tiene \emph{solo texto}. Suele combinar \textbf{tres elementos:} \textbf{(1) texto en párrafos} para narrar, explicar, argumentar; \textbf{(2) tablas} para mostrar datos comparativos o relacionados; \textbf{(3) imágenes} para ilustrar, ejemplificar, hacer ameno. Insertar una tabla en Word es muy rápido: \texttt{Insertar} → \texttt{Tabla} → elegir filas y columnas. Insertar una imagen también: \texttt{Insertar} → \texttt{Imágenes} → escoger desde tu computador. Lo que diferencia a un usuario \emph{básico} de uno \emph{intermedio} es saber \textbf{ajustar} estos elementos: que la tabla se vea limpia, que la imagen tenga texto fluyendo a su lado en vez de bloquear todo. Hoy aprendes esas técnicas.
\end{softbox}

\begin{softbox}{milcMostaza}{milcGris}{Pregunta-puente}
\textit{Si tu profe te pide entregar un trabajo sobre \emph{``los 5 ríos más importantes de Colombia''}, ¿\textbf{cómo lo presentas mejor}? ¿En 2 párrafos de texto seguido, o en una tabla con columnas \emph{nombre · longitud · ubicación · datos curiosos}? ¿Por qué?}
\end{softbox}

\begin{softbox}{milcVerde}{milcGris}{Saber y saber hacer}
Al terminar la clase: (1) podrás \textbf{identificar} cuándo usar tabla y cuándo imagen; (2) sabrás \textbf{aplicar} la inserción de ambos elementos en Word/Docs; (3) podrás \textbf{evaluar} las 5 reglas de uso de imágenes; (4) habrás \textbf{creado} un documento con tabla, imágenes y texto integrado.
\end{softbox}

\small
\begin{tabularx}{\textwidth}{>{\bfseries}p{3.0cm}Y}
\rowcolor{milcGradePrimary!12}Autor & Dr. Álvaro Cárdenas Orozco\\
DBA & Produce contenidos digitales y valida información de fuentes web (MEN, T\&I 6°)\\
Referentes & Saber ancestral de la maestra rural · Lectura crítica · Soberanía informacional\\
Producto final & Un \textbf{documento} con: \textbf{(a) una tabla} de 5 filas × 4 columnas que muestre datos reales (ejemplo: \emph{horario semanal del colegio} con materia, día, hora, salón); \textbf{(b) 2 imágenes} insertadas con \emph{ajuste de texto alrededor} (una a la izquierda, otra a la derecha), cada una con \emph{pie de imagen} debajo; \textbf{(c) 2 párrafos de texto} relacionados con las imágenes que fluyen alrededor de ellas. Más \textbf{en el cuaderno}: dibujo del documento + 5 reglas de uso de imágenes que aprendiste.\\
\end{tabularx}
\normalsize

\puente{Hoy haces 4 cosas: decides cuándo conviene tabla vs imagen, aprendes a insertar y ajustar ambos, aplicas el ajuste de texto alrededor, y armas tu documento mixto.}

\begin{titlebox}{milcGradeDark}
Ruta MILC: así vamos a aprender
\end{titlebox}
\begin{tabularx}{\textwidth}{>{\centering\arraybackslash}X>{\centering\arraybackslash}X>{\centering\arraybackslash}X>{\centering\arraybackslash}X}
\phasecard{1}{milcVerde}{milcGris}{Escuta\\[-1mm]\small Observo, escucho y pregunto.} &
\phasecard{2}{milcTurquesa}{milcGris}{Sistematización\\[-1mm]\small Pensamiento computacional.} &
\phasecard{3}{milcMagenta}{milcGris}{Praxis\\[-1mm]\small Construyo y aplico.} &
\phasecard{4}{milcVino}{milcGris}{Evaluación\\[-1mm]\small Reflexión liberadora.}
\end{tabularx}


\puente{Empezamos con 6 situaciones para decidir: ¿texto, tabla o imagen?}

\begin{titlebox}{milcVerde}
Fase 1 · Escuta
\end{titlebox}

\begin{softbox}{milcVerde}{milcGris}{Escena para escuchar}
\textbf{Actividad 1 · IDENTIFICA --- 6 situaciones para clasificar} (10 min · individual). En el cuaderno, dibuja 3 columnas: \emph{TEXTO}, \emph{TABLA}, \emph{IMAGEN}. Lee estas 6 situaciones y clasifica cada una en la columna que mejor corresponde. \textbf{(1)} Explicar las 4 estaciones del año con sus características. \textbf{(2)} Mostrar paso a paso cómo armar un caballero de Lego. \textbf{(3)} Comparar 5 países por población, capital, idioma y moneda. \textbf{(4)} Contar la historia de tu fin de semana. \textbf{(5)} Mostrar cómo se ve por dentro un volcán. \textbf{(6)} Listar los presidentes de Colombia con año de mandato. \emph{Tip:} algunas podrían ir en 2 columnas; escoge la PRINCIPAL.
\end{softbox}

\checkbox Dibujé las 3 columnas en el cuaderno.\\
\checkbox Clasifiqué las 6 situaciones, una por columna.\\
\checkbox Pensé cuál es la herramienta principal para cada caso.

\begin{infoband}{milcVerde}
\textbf{Lo que va al cuaderno (Actividad 1 · IDENTIFICA).} Título: «Actividad 1 --- 6 situaciones: ¿texto, tabla o imagen?». \textbf{Formato:} tabla de 3 columnas con las 6 situaciones distribuidas. \textbf{Extensión:} 6 situaciones. \textbf{Sabes que terminaste cuando} puedes justificar cada decisión.
\end{infoband}

\puente{Después aprendes a insertar tablas y a manejar las imágenes con ajuste de texto.}

\begin{titlebox}{milcTurquesa}
Fase 2 · Sistematización con pensamiento computacional
\end{titlebox}

\begin{softbox}{milcTurquesa}{milcGris}{Cuatro pilares aplicados al tema}
Ya intuiste cuándo conviene cada elemento. Ahora aprendes las \textbf{reglas técnicas} para insertarlos y ajustarlos bien.
\end{softbox}

\small
\begin{tabularx}{\textwidth}{>{\bfseries\RaggedRight\arraybackslash}p{3.6cm}Y}
\rowcolor{milcTurquesa!90}\color{white}Pilar & \color{white}\bfseries Aplicación al tema\\
1. Descomposición & \textbf{Cuándo usar TABLA.} Cuando tienes \emph{datos comparables} en categorías. Ejemplo clásico: \emph{lista de 5 ríos con 4 datos cada uno = 5 filas × 4 columnas}. \textbf{En Word:} \texttt{Insertar → Tabla → elegir tamaño}. Aparece la cuadrícula. Llenas las celdas. La fila de encabezado se pone en \textbf{negrita} y con \emph{fondo gris claro}. \textbf{Cuándo usar IMAGEN.} Cuando algo se ve mejor que se explica: un mapa, un esquema, una foto que ilustra, un diagrama. \textbf{En Word:} \texttt{Insertar → Imágenes → seleccionar archivo}. Por defecto la imagen llega \emph{enorme} y desplaza todo el texto.\\
2. Reconocimiento de patrones & \textbf{Ajuste de texto alrededor de la imagen.} El truco profesional: \textbf{haz que el texto fluya al lado de la imagen} en vez de empujarla. En Word: clic en la imagen → aparece el ícono \emph{Opciones de diseño} (una arquetilla) → escoge \texttt{Cuadrado} (texto rodea la imagen) o \texttt{Estrecho} (más pegado). Después arrastras la imagen donde la quieres. Tamaño ajustado: clic en imagen → arrastrar de las esquinas (no de los lados, deforma).\\
3. Abstracción & \textbf{Las 5 reglas de uso de imágenes.} (1) \textbf{Resolución suficiente:} no uses imágenes pequeñas que se ven borrosas al ampliar. (2) \textbf{Pie de imagen:} debajo de la imagen, escribe en \emph{tamaño 10, cursiva}, una línea que diga qué muestra (\emph{``Figura 1: mapa de Colombia''}). (3) \textbf{Derechos de autor:} no uses cualquier imagen de Google. Usa \emph{Unsplash}, \emph{Pixabay}, \emph{Pexels} (gratuitas y de uso libre). (4) \textbf{Tamaño moderado:} máximo 1/3 del ancho de la página, dejando espacio para texto al lado. (5) \textbf{Relevancia:} cada imagen debe \emph{aportar algo}; si es decoración pura, mejor no la pongas.\\
4. Algoritmo conceptual & \textbf{Errores frecuentes que se evitan.} (1) \emph{Insertar imagen y dejarla bloqueando todo el ancho} cuando solo necesitabas 1/3 de espacio. (2) \emph{Tabla sin encabezado en negrita} --- los lectores no saben qué es cada columna. (3) \emph{Imagen sin pie} --- después no se sabe qué muestra. (4) \emph{Tabla con muchas celdas vacías} --- mejor usar otra cosa. (5) \emph{Imagen con marca de agua o copyright} --- se ve poco profesional y puede ser ilegal usarla.\\
\end{tabularx}
\normalsize

\begin{drawbox}{milcTurquesa}{Cuándo + cómo + 5 reglas}
\textbf{TABLA:} datos comparables. Insertar → Tabla → elegir filas y columnas. Encabezado en negrita. \\[1mm]
\textbf{IMAGEN:} cuando se ve mejor que se explica. Insertar → Imágenes. Ajuste de texto: Cuadrado o Estrecho. \\[1mm]
\textbf{5 reglas:} resolución · pie · derechos · tamaño moderado · relevancia.
\end{drawbox}

\begin{softbox}{milcMostaza}{milcGris}{Errores comunes y su corrección}
\textbf{Mitos sobre tablas e imágenes.} (1) \emph{``Más imágenes = más profesional''} --- Falso. Las imágenes irrelevantes restan. (2) \emph{``Las imágenes de Google las puedo usar libres''} --- Falso. Muchas tienen copyright. Usa Pixabay, Pexels, Unsplash. (3) \emph{``Una tabla con 2 filas se justifica''} --- Falso. Mínimo 3 filas. Con 2 mejor 2 párrafos. (4) \emph{``El texto envuelve solo, no hay que configurarlo''} --- Falso. Por defecto la imagen empuja todo. Hay que cambiar el ajuste a \emph{Cuadrado}. (5) \emph{``Las imágenes grandes se ven mejor''} --- Falso. Las gigantes parecen amateur. 1/3 del ancho es elegante.
\end{softbox}

\begin{infoband}{milcTurquesa}
\textbf{Lo que va al cuaderno (Actividad 2 · EXPLICA).} Título: «Actividad 2 --- 5 reglas de imágenes + cómo se inserta tabla». \textbf{Formato:} bloque A con las 5 reglas numeradas (1 línea cada una); bloque B con el camino para insertar tabla (Insertar → Tabla → escoger) y para insertar imagen (Insertar → Imágenes). \textbf{Extensión:} 7-8 líneas. \textbf{Sabes que terminaste cuando} podrías hacer un tutorial al compañero de al lado.
\end{infoband}

\puente{Con todo claro, creas tu documento con los 3 elementos integrados.}

\begin{titlebox}{milcMagenta}
Fase 3 · Praxis con pensamiento computacional
\end{titlebox}

\begin{softbox}{milcMagenta}{milcGris}{Construyendo el producto con los cuatro pilares}
\textbf{Actividad 3 · CREA --- Documento mixto con tabla + imágenes + texto} (32 min · individual). Vas a crear un documento que demuestra los 3 elementos integrados. \emph{Tema sugerido}: \emph{``Mi colegio''} o \emph{``Mi cuadra''}. Tú escoges.
\end{softbox}

\small
\begin{tabularx}{\textwidth}{>{\bfseries\RaggedRight\arraybackslash}p{3.6cm}Y}
\rowcolor{milcMagenta!90}\color{white}Pilar & \color{white}\bfseries Aplicación al producto\\
Descomposición & \textbf{Paso 1 --- Documento nuevo + título.} Abre Word o Docs. Título centrado en negrita azul oscuro tamaño 18: \emph{``Mi colegio --- texto, tabla e imágenes''}.\\
Patrones & \textbf{Paso 2 --- Inserta la tabla del horario.} Escribe un subtítulo: \textbf{``Mi horario semanal''}. Debajo, inserta una tabla de \textbf{5 filas × 4 columnas} (\texttt{Insertar → Tabla → 5 × 4}). Encabezado de columnas: \emph{Materia | Día | Hora | Salón}. Llena las 4 filas restantes con 4 materias reales tuyas y sus datos. Pon el encabezado en \textbf{negrita}.\\
Abstracción & \textbf{Paso 3 --- Inserta primera imagen + pie.} Salta una línea. Escribe un subtítulo: \textbf{``Vista del colegio''}. Inserta una imagen (\texttt{Insertar → Imágenes}). Si no tienes foto del colegio, busca en \emph{Pixabay} una imagen genérica de \emph{escuela} (gratis). Selecciónala. Clic en ella → \emph{Opciones de diseño} → \textbf{Cuadrado}. Arrastra la imagen al \emph{lado izquierdo} y ajusta tamaño a 1/3 del ancho. Debajo de la imagen, en cursiva tamaño 10: \emph{``Figura 1: el colegio desde el patio.''}.\\
Algoritmo de armado & \textbf{Paso 4 --- Escribe 2 párrafos junto a la primera imagen.} A la derecha de la imagen, escribe 2 párrafos cortos (3-4 líneas cada uno) sobre tu colegio: en qué barrio queda, cuántos estudiantes hay, qué te gusta. El texto debe \emph{fluir alrededor} de la imagen (eso se logra con el ajuste \emph{Cuadrado} que aplicaste).\\
\end{tabularx}
\normalsize

\begin{drawbox}{milcMagenta}{Antes de mostrar tu documento mixto}
\checkbox Tiene título principal formateado. \\[1mm]
\checkbox La tabla 5x4 está completa con encabezado en negrita. \\[1mm]
\checkbox Hay 2 imágenes con ajuste cuadrado (texto fluye al lado). \\[1mm]
\checkbox Cada imagen tiene pie en cursiva tamaño 10. \\[1mm]
\checkbox Las imágenes son de fuente libre (Pixabay/Pexels/Unsplash). \\[1mm]
\checkbox Está guardado con nombre correcto en la carpeta del periodo.
\end{drawbox}

\begin{softbox}{milcMagenta}{milcGris}{Plantilla de trabajo}
\textbf{Estructura del documento.} \textit{Título centrado.} \textit{Subtítulo + tabla 5×4 del horario.} \textit{Imagen izquierda + pie + 2 párrafos a la derecha.} \textit{Imagen derecha + pie + 1 párrafo a la izquierda.} \textit{Guardado:} 2026-mi-colegio-mixto.docx.
\end{softbox}

\begin{infoband}{milcMagenta}
\textbf{Lo que va al cuaderno (Actividad 3 · CREA).} Título: «Actividad 3 --- Mi documento mixto: tabla + imágenes». \textbf{Formato:} boceto del layout + datos del horario + las 5 reglas de imágenes anotadas. \textbf{Extensión:} 1 página. \textbf{Sabes que terminaste cuando} el documento muestra integración real (no son elementos sueltos).
\end{infoband}

\puente{El producto es ese documento + las 5 reglas de imágenes en el cuaderno.}

\begin{titlebox}{milcMostaza}
Producto de la sesión
\end{titlebox}
\textbf{Documento mixto con tabla 5x4 + 2 imágenes con texto fluyendo}.
\begin{softbox}{milcMostaza}{milcGris}{Criterios de calidad}
(1) La tabla tiene 5 filas, 4 columnas y encabezado en negrita. (2) Hay 2 imágenes con ajuste cuadrado (texto alrededor). (3) Cada imagen tiene pie en cursiva tamaño 10. (4) Las imágenes son de fuente libre (no copyright). (5) El documento está guardado con nombre correcto.
\end{softbox}

\puente{Antes de salir, verificas que la tabla y las imágenes están bien insertadas con el texto fluyendo.}

\begin{titlebox}{milcVino}
Fase 4 · Evaluación liberadora — cinco dimensiones
\end{titlebox}

\begin{softbox}{milcVino}{milcGris}{Sentido de la evaluación}
Evaluar no es solo revisar una nota. Es reconocer qué aprendí, qué cambió en mí, qué emociones manejé, cómo me relaciono con la ciudad y la comunidad, y qué le dejaré a quien viene detrás.
\end{softbox}

\textbf{Mi reflexión liberadora en cinco dimensiones.} Completa cada frase iniciada:

\small
\begin{tabularx}{\textwidth}{>{\bfseries\RaggedRight\arraybackslash}p{3.4cm}Y>{\RaggedRight\arraybackslash}p{4.6cm}}
\rowcolor{milcVino}\color{white}Dimensión & \color{white}\bfseries Mi respuesta (frame starter) & \color{white}\bfseries Algo que descubrí\\
1. Personal & Lo que más cambió en mí fue \linead{6.5cm} & \linead{4.2cm}\\
2. Emocional & La emoción más fuerte fue \linead{2.5cm} y la regulé \linead{3cm} & \linead{4.2cm}\\
3. Ciudadana & En democracia esto importa porque \linead{6cm} & \linead{4.2cm}\\
4. Local (Cartago) & En mi barrio sería útil para \linead{6.5cm} & \linead{4.2cm}\\
5. Intergeneracional & A mi abuela/abuelo le diría \linead{6.5cm} & \linead{4.2cm}\\
\end{tabularx}
\normalsize

\begin{infoband}{milcVino}
\textbf{Para la próxima sustentación.} Vuelve a esta tabla en seis meses. Verás que las respuestas cambian — no porque lo aprendido cambie, sino porque tú cambias. La evaluación liberadora es una conversación contigo a través del tiempo.
\end{infoband}


\puente{Cerramos con 3 voces que te ayudan a entender por qué mostrar es a veces más que decir.}

\begin{titlebox}{milcGradeDark}
Triángulo de pensamiento · cierre filosófico
\end{titlebox}

\begin{softbox}{milcVino}{milcGris}{Dussel · lente del nosotros}
\textit{````Hay verdades que se dicen con palabras, otras se muestran con imágenes, otras se ordenan en filas. Saber cuándo usar cada una es saber comunicar.''''}

Doña Mercedes sabía que enseñar a una niña qué es un río era distinto a enseñarle qué es una vaca: \emph{un río se muestra con mapa; una vaca con foto; los ríos del país en una tabla}. La sabiduría está en \textbf{saber qué herramienta para qué verdad}. Hoy con procesadores tienes todas a tu disposición. Saber cuándo usar tabla, cuándo imagen, cuándo solo texto, eso es \emph{maestría comunicativa}.

\textbf{¿Estoy usando la herramienta correcta para cada idea que quiero comunicar?}
\end{softbox}
\linead{\linewidth}\\[1mm]
\linead{\linewidth}

\begin{softbox}{milcMostaza}{milcGris}{Estoicismo · Marco Aurelio (emperador que documentaba con orden) · cuidado interior}
\textit{````Una tabla bien hecha vale lo que dos páginas de palabras. Eso es economía de pensamiento.''''}

Hay gente que escribe largo cuando podría escribir corto. \textbf{La tabla obliga a la economía:} cabeceras claras, datos justos, sin relleno. \emph{Pensar en tablas te entrena en pensar conciso.} En el futuro, cuando escribas informes, ensayos, reportes, esta habilidad te ahorra horas a ti y a quienes te leen.

\textbf{¿Qué cosa que escribí en párrafo se entendería mejor en tabla?}
\end{softbox}
\linead{\linewidth}\\[1mm]
\linead{\linewidth}

\begin{softbox}{milcTurquesa}{milcGris}{Floridi · lente de la infoesfera}
\textit{````La imagen sin contexto es ruido. La tabla sin encabezado es caos. La integración es lo que hace al documento profesional.''''}

Una imagen suelta sin pie de foto NO comunica: el lector se pregunta \emph{``¿qué es esto?''}. Una tabla sin encabezado tampoco: \emph{``¿qué significa esta columna?''}. \textbf{Lo profesional es integrar:} texto que conecta, tabla con cabeceras, imagen con pie. Cuando todo está integrado, el documento se lee como pieza única, no como collage.

\textbf{Cuando reviso mis documentos, ¿está cada elemento integrado o son piezas sueltas?}
\end{softbox}
\linead{\linewidth}\\[1mm]
\linead{\linewidth}

\begin{titlebox}{milcVino}
Compromiso final
\end{titlebox}

\small
\begin{tabularx}{\textwidth}{>{\RaggedRight\arraybackslash}p{3.0cm}YYY}
\rowcolor{milcVino}\color{white}\bfseries Criterio & \color{white}\bfseries Lo logré & \color{white}\bfseries En proceso & \color{white}\bfseries Necesito apoyo\\
Comprensión del tema & Explico ideas centrales con mis palabras. & Reconozco ideas pero falta precisión. & Necesito repasar conceptos base.\\
Pensamiento computacional & Apliqué los 4 pilares al diseñar mi producto. & Apliqué algunos pilares. & Necesito releer la fase 2.\\
Aplicación y producto & Mi producto cumple los criterios y se puede revisar. & Producto funcional con ajustes pendientes. & Aún en construcción.\\
Reflexión 5 dimensiones & Respondí cada dimensión con honestidad. & Respondí algunas con detalle. & Mis respuestas son muy generales.\\
Triángulo de pensamiento & Conecté las 3 voces con mi trabajo real. & Conecté al menos una. & No relaciono las voces con mi proceso.\\
\end{tabularx}
\normalsize

\textbf{Hoy entendí que:}\\[1mm]
\linead{\linewidth}\\[1mm]
\linead{\linewidth}

\textbf{Mi compromiso para la próxima sesión es:}\\[1mm]
\linead{\linewidth}\\[1mm]
\linead{\linewidth}

\begin{softbox}{milcMostaza}{milcGris}{Cita de despedida · Marco Aurelio}
\textit{``El obstáculo en el camino se vuelve el camino. Lo que estorba al camino se vuelve camino.''}\\[1mm]
Cada error de hoy, bien observado, se convierte en pieza del camino que pisarás mañana.
\end{softbox}

\small
\begin{tabularx}{\textwidth}{>{\bfseries}p{3.6cm}Y}
\rowcolor{milcGradeDark!12}Nombre completo: & \linead{10cm}\\
Fecha: & \linead{4cm} \quad Curso/grupo: \linead{4cm}\\
Mi firma: & \linead{9cm}\\
\end{tabularx}
\normalsize

\begin{infoband}{milcGradeDark}
\textbf{Para la próxima sesión.} Esta semana, en cualquier trabajo digital, considera si una tabla o una imagen mejoraría la entrega. La próxima clase aprendes \textbf{encabezados, pies de página y numeración}: los toques finales que hacen un documento profesional --- como las solapas y la portada de un libro.
\end{infoband}

\end{document}
